\section{Review Limitations}

Five limitations bound the reading of the results and qualify the reach of the conclusions.

\noindent\textbf{Extraction granularity.} The thematic synthesis derives from the structured fields of the records (artifacts, roles, execution, validation, portability, risks, nature, and evidence) and from a consolidated narrative finding per study, produced in a second extraction pass. Twelve studies have a finding drawn from a substantive inclusion justification, and the other twenty-five receive a finding consolidated from the structured fields, transparently flagged as such. Where the first-pass justification is generic, the reading still rests in part on the known public identity of well-documented frameworks, such as ChatDev, MetaGPT, AutoCodeRover, Reversa, and AgileGen. A full-text, field-by-field extraction for the remaining twenty-five studies would deepen the synthesis further.

\noindent\textbf{Quality assessment by category, not by full text.} We applied a four-item source-category quality assessment (Table~\ref{tab:qualidade}) that weights the strength of the evidence and flags the lower-confidence studies, but it does not replace a formal risk-of-bias instrument applied to the full text of each study. The ten arXiv preprints not yet peer-reviewed and the two studies with medium extraction confidence enter the synthesis with a cautious reading, and no theme rests exclusively on them; even so, a full-text reappraisal with a validated instrument would refine the weighting.

\noindent\textbf{Language and database scope, and publication bias.} The protocol restricts the review to English and Portuguese, in the 2018 to 2026 window. The direct arXiv search failed in part due to access instability, with preprint coverage recovered via Semantic Scholar and the arXiv DOI, which may have left out very recent preprints not indexed by those routes. A publication bias toward frameworks with positive results is likely and measurable in the corpus itself: 26 of the 37 studies (70\%) claim empirical evaluation, almost always favorable to the proposed framework, while studies of failure, abandonment, or non-adoption are practically absent; only the adoption axis of theme T5, with six studies, problematizes organizational integration. Widening the search to additional databases and lifting the language restriction is future work that could balance this asymmetry.

\noindent\textbf{Operational pending item outside the corpus.} One source approved at screening lacked full text obtained at the cutoff date and was kept as an operational pending item, outside the extraction and this synthesis. The corpus effectively synthesized is 37 studies; incorporating that source later would adjust marginal counts without changing the thematic structure.

\noindent\textbf{Partial snowballing saturation.} Snowballing was stopped at the protocol limit of three rounds. The saturation curve (Figure~\ref{fig:saturacao}) shows decreasing inclusions, from 25 in the initial search to 6, 3, and 3 across the three rounds: a clear but nonzero deceleration. The last rounds still added studies, which indicates partial rather than full saturation; additional rounds could surface relevant studies, especially outside the communities of origin of the seeds.
